\chapter{Diseño}

La mayoría del diseño se encuentra en diagramas, en los directorios de documentación correspondientes.

\section{Decisiones fundamentales}

Tras realizar el análisis, se consideraron las siguientes decisiones como las más adecuadas y que mejor se apegan a este proyecto.

\begin{itemize}
	\item \textbf{Lenguaje de programación:} El cliente y el servidor se programarán en \texttt{C\#} con el framework de \texttt{.NET}, ya que el lenguaje es orientado a objetos y el framework tiene bibliotecas útiles para el protocolo TCP/IP, serialización y deserialización de JSON, entre otras ventajas.
	\item \textbf{Programas:}
	\begin{itemize}
		\item \textbf{Servidor:} Se utilizará el patrón de diseño MVC y no se contará con interfaz gráfica, sino con salida de texto a consola.
		\item \textbf{Cliente:} Se utilizará el patrón de diseño MVVM y contará con interfaz gráfica, para lo cual se utilizará el framework de \texttt{Avalonia}.
	\end{itemize}
	\item \textbf{Sistema de construcción:}
	\begin{itemize}
		\item \textbf{Compilación:} \textit{MSBuild} a través del framework de \texttt{.NET}.
		\item \textbf{Generación de documentación:} Herramienta \textit{Doxygen}.
		\item \textbf{Pruebas unitarias:} \textit{MSTest}, integrado en el framework de \texttt{.NET}.

	\end{itemize}
	\item \textbf{Sistema de control de versiones:} \textit{Git} y \textit{GitHub} para repositorio remoto.
\end{itemize}


\section{Directorios}

\dirtree{%
    .1 FC\_MYP\_Proyecto1/.
        .2 Docs/.
            .3 Process/.
            .3 AnalysisDesign/.
            .3 Report/.
            .3 CodeDocs/.
                .4 Client/.
                .4 Server/.
        .2 Client/.
            .3 Model/.
            .3 ViewModel/.
            .3 View/.
        .2 Server/.
            .3 Model/.
            .3 Controller/.
            .3 View/.
}
\vspace{2mm}

A continuación, una explicación general del contenido de cada directorio:

\begin{itemize}
	\item \textbf{FC\_MYP\_Proyecto1:} Destinado a contener todos los directorios del proyecto y archivos de configuración (\texttt{.gitignore}, \texttt{doxygen}, etc.)
	\item \textbf{Docs:} Guarda la documentación de todo el proceso llevado a cabo, desde el análisis hasta la distribución. 
	\begin{itemize}
		\item \textbf{Process:} Documenta el proceso de cómo se desarrolló el proyecto y qué flujos de trabajo se siguieron.
		\item \textbf{AnalysisDesign:} Elementos necesarios para el análisis y diseño.
		\item \textbf{Report:} Contiene este reporte y toda su organización.
		\item \textbf{CodeDocs:} Contiene la documentación del código, tanto del cliente como del servidor, generada por la herramienta \textit{Doxygen}.
	\end{itemize}
	\item \textbf{Client:} Contiene el ejecutable del cliente y todos los directorios necesarios para su implementación.
	\item \textbf{Server:} Contiene el ejecutable del servidor y todos los directorios necesarios para su implementación.
\end{itemize}

\section{Objetos}

\subsection{Controlador del servidor}
Como ya hemos discutido, seguiremos un flujo 
\begin{center}
	recibir mensaje $\to$ decodificarlo $\to$ traducirlo a una acción-respuesta $\to$ ejecutar la acción-respuesta
\end{center}
 entonces el controlador va a ser el encargado de cumplir con esta tarea, además de proveer eventos para notificar a la vista y que se puedan ver los mensajes recibidos y enviados desde la consola del equipo donde se esté corriendo el servidor. Para que pueda comunicar la recepción y envío de mensajes las componentes que use también deben tener una manera de notificárselo.
 \\
 
 Como el servidor puede correr y parar, eso hará el controlador. Al correr abre la comunicación, indica que está activo y se pone a aceptar conexiones de manera asíncrona. Una vez acepta una conexión, maneja al cliente, lo cual significa que dado que la conexión se estableció, ahora tiene que escuchar sus mensajes, lo cual realiza en un bucle que existe mientras la conexión exista, igualmente de manera asíncrona, para no quedarse estancado en un sólo cliente. Cuando el cliente se desconecta, por él mismo o porque el servidor lo desconectó, el bucle termina y deja de ser manejado.
\\

\textbf{Pseudocódigo}

\begin{verbatim}
	class ServerController:
	  event MessageReceived(msg : string, from : string)
	  event MessageSent(msg : string, to : string)
	  
	  field tcpl : TcpListener
	  field isActive : bool
	  field encoder : IEncoder<string>
	  field translator : Translator
	  
	  constructor ServerController(ip : string, port : string):
	    tcpl = Crear TcpListener con IP ip conectada al puerto p.
	    isActive = false.
	    encoder = Crear implementación de IEncoder<string>
	    para bytes<-> UTF-8.
	    translator =  nuevo Translator.
	    
	  method Run():
	    StartListening()
	    isActive = true.
	    Mientras isActive sea true:
	      socketCliente = esperar a recibir conexión de tcpl de forma asíncrona
	      como TcpClient.
	      cliente = Crear un Client con nombre vacío, socket socketCliente
	      y subirlo a la bd de clientes conectados. Marcarlo como conectado.
	      manejarClienteTarea = HandleClient(cliente)
	    Desconectar a todos los clientes
	      
	  method Stop():
	    isActive = false.
	  method StartListening():
	    Comenzar el tcpl.
	  
	  method StartListening():
	    comenzar tcpl.
	    
	  method HandleClient(c : TcpClient):
	    bytes = ReceiveData().
	    msg = Decodificar bytes con el encoder.
	    strategy = Traducir msg con el translator.
	    Ejecutar strategy
	    
	  método Receivedata(c : Client)
	    stream = Obtener el stream del socket de c.
	    datos codificados = Obtener los datos del stream (los primeros n bytes)
	    regresar datos codificados.  
\end{verbatim}

\subsection{Codificadores}

Consiste en una interfaz \texttt{IEncoder<T>} que provee métodos para codificar y decodificar un mensaje, de algún tipo a bytes y viceversa, respectivamente, y en un codificador \texttt{UTF8Enconder} que codifica de \texttt{byte[]} a \texttt{string} en formato UTF-8.
\\

\textbf{Pseudocódigo}
\begin{verbatim}
	clase IEncoder<T>:
	  
	  método decode(data : byte[]) : T.
	     
	  método encode(data : T) : byte[].
	    
\end{verbatim}

\subsection{Estrategias de acción respuesta}
Consiste en una interfaz \texttt{IARStrategy} y un conjunto de doce estrategias (implementaciones de dicha interfaz) que realizan las operaciones relacionadas a los tipos de mensajes que puede recibir el servidor. Supongamos que ya identificamos la intención del mensaje de un cliente recibido por el servidor, sabemos de qué cliente vino y conocemos los parámetros del mensaje (el \texttt{username} en caso de identificar un cliente, el \texttt{status} en caso de cambiarle el estado, etc.). Entonces el problema se reduce a realizar consultas y modificaciones en la base de datos temporal de nuestro servidor y en responderle al cliente que envió el mensaje y notificarle al resto sobre la operación realizada, en caso de ser requerido. Dado que conocemos al cliente, podemos enviarle el mensaje, y podemos acceder a la base de datos mediante el DAO correspondiente, pero eso le corresponde a las estrategias de acción-respuesta concretas.
\\

\textbf{Pseudocódigo}
\begin{verbatim}
	interface IARStrategy:
		method Execute()
\end{verbatim}
Para las doce estrategias concretas no hay pseudocódigo, ya que se deduce del diagrama de flujo del servidor, el cual puede encontrarse en el repositorio en la ruta \texttt{Docs/Analysis\_Design/ServerFlow/server\_flow\_diagram.drawio}.

\subsection{Traductor de mensajes}
Consiste en un objeto \texttt{MessageTranslator} que traduce un \texttt{msg} (\texttt{string} que es el mensaje recibido por el cliente) a una estrategia de acción respuesta \texttt{IARStrategy} a la cual inicializa con sus parámetros requeridos (el \texttt{username} para \texttt{IdentifyUserSt}, el \texttt{status} para \texttt{ChangeStatusSt}, etc.) a través del contenido adicional del mensaje. Lo lee en formato JSON, por lo cual se encarga de deserializarlo, y una vez identificando el tipo de operación lo vuelve a deserializar en un contenedor de datos del tipo de respuesta concreta que está recibiendo, cuyos campos son los argumentos que se le pasarán a la estrategia concreta.
\\

\textbf{Pseudocódigo}

\begin{verbatim}
	class MessageTranslator:	
	  method Translate(msg : string, c : Client) : IARStrategy:
	    mensajeBase = Deserializar msg.
	    estrategia = null;
	    Extraer el campo `type' del mensajeBase y comparar:
	      Si es ``IDENTIFY'':
	        contenedorMsgIdentify = Deserializar msg.
	        estrategia = Crear IdentifySt pasándole el campo ``username''
	        de mensajeIdentificar y al cliente c.
	      .
	      .
	      .
	      Si es ``DISCONNECT'':
	        estrategia = Crear DisconnectSt pasándole al cliente c.
	      Si no es ninguno:
	        Excepción ``Mensaje no traducible''
	    regresar estrategia.
\end{verbatim}


\subsection{Contenedores de datos de mensaje}
Consiste en una clase abstracta \texttt{MessageDataBase} y un conjunto de, a lo más, doce clases que la heredan que especifican campos del tipo de mensaje proveniente del cliente, por cada uno de los doce tipos. Sirven como recurso para almacenar datos deserializados de un mensaje recibido por el servidor.
\\

\textbf{Pseudocódigo}
\begin{verbatim}
	abstract class MessageDataBase:
	  field type : string.
	  Property Type => type.
	  
	  constructor MessageDataBase(t : string):
	    type = t.  
\end{verbatim}

Las clases concretas que la heredan sólo aumentan la cantidad de campos que pueden ser serializados, dependiendo del tipo de mensaje que se espera.

\subsection{Base de datos y DAOs}
Consisten en un singleton \texttt{ServerDB} y un conjunto de DAOs (\texttt{ConnectionsDAO}, \texttt{IdentifiedDAO}) que permiten el acceso a dicha ``base de datos'' temporal, la cual almacena usuarios conectados, identificados, salas públicas y privadas (chats entre dos personas). Sus métodos se especifican en el diagrama de clases del servidor, en la misma ruta que el diagrama de flujo de las estrategias.

\subsection{Servidor}

 Al final, nuestra vista del servidor se ve de la siguiente forma:
 
\begin{verbatim}
	clase Servidor:
	    método Main():
	      Asignar comando ``ctrl + c'' para finalizar el servidor.
	      puerto = Preguntar el puerto donde se ejecutará el servidor.
	      controlador = Crear el ServerController(IP local, puerto).
	      Correr el controlador.
	      Imprimir mensaje de despedida.
\end{verbatim}

